% !PRINTED 331
\weberpage{331}
\begin{center}
\S.\,89.

\gesperrt{Inflexion points of a curve of the third order.}
\end{center}

Let $f=0$ now be the equation of a curve of the third order
without a singular point. We place two of the vertices of the
coordinate triangle $\xi_1, \xi_2$ at two of the nine inflexion points,
and take the corresponding inflexion tangents for the
sides $x_2, x_1$. Then, if $x_1=0$ or $x_2=0$ is set, $f$ must
reduce to the term $h\,x_3^3$, in which $h$ is a constant,
and consequently all terms that have a coefficient different
from zero in the cubic form, excepting
$h\,x_3^3$, must be divisible by $x_1\,x_2$. Accordingly $f$ takes the form

\weq{1}{f\,(x_1,x_2,x_3) = x_1\,x_2\,(a_1\,x_1+a_2\,x_2+a_3\,x_3) + h\,x_3^3,}

in which $a_1, a_2, a_3$ are constants. From this, however, it follows that
the straight line

\[
a_1\,x_1 + a_1\,x_2 + a_3\,x_3 = 0
\] % [sic] printed coefficient of x_2 is a_1; expected a_2

is likewise an inflexion tangent of the curve $f$, and that its point of
intersection with the line $x_3=0$ is a third inflexion point.
This point is different from the first two $\xi_1, \xi_2$, because
neither $a_1$ nor $a_2 = 0$ can hold if the curve has no singular
point. We shall express this important result in the form
of a theorem:

1. \gesperrt{If one joins any two inflexion points of a
curve of the third order without a singular point
by a straight line, then this straight line cuts
the curve in a third inflexion point.}

Since there are nine inflexion points, four straight lines
issue from each of these points, each of which contains two
further inflexion points. There are nine such pencils
of four lines, and because each of these lines occurs in three pencils,
the number of the lines is $9\,.\,4 : 3 = 12$. We
can accordingly supplement Theorem 1. as follows:

2. \gesperrt{The nine inflexion points of a curve of the third
order lie by threes on twelve straight
lines, and through each inflexion point pass
four of these lines.}
